% !PRINTED 553
\weberpage{553}

Because of (4) this theorem need only be proved for
$\mu_n = n$, since the series $\Sigma\, \varepsilon_n\theta/n^{s+\frac{1}{2}}$ converges unconditionally and
therefore has this property by known elementary theorems.

Let us therefore set
\weq{7}{\sigma = \frac{\varepsilon_1}{1^s} + \frac{\varepsilon_2}{2^s} + \frac{\varepsilon_3}{2^s} + \cdots} % [sic] third denominator printed 2^s, not 3^s
and, by (5), $\varepsilon_n = \gamma_n\sqrt{n} - \gamma_{n-1}\sqrt{n-1}$; then
\weq{8}{\sigma = \frac{\varepsilon_1}{1^s} + \gamma_2\sqrt{2}\left(\frac{1}{2} - \frac{1}{3^s}\right) + \gamma_3\sqrt{3}\left(\frac{1}{3^s} - \frac{1}{4^s}\right) + \cdots} % [sic] German source prints $2$, not $2^s$, as the first denominator in the parenthesis

Since now, for an infinitely large $n$,
\[
\sqrt{n}\left(\frac{1}{n^s} - \frac{1}{(n+1)^s}\right) = \frac{s}{n^{s+\frac{1}{2}}}
\]
holds, the series (8) is unconditionally convergent as long as $s > \frac{1}{2}$,
and from this our assertion follows for this series also
by the same elementary theorems.

\begin{center}\textbf{\S\,147. Development of $f, f_1, f_2$.}\end{center}

The series (12), § 145 are now in this case. Let us consider,
for example, the second of them, from which, as we have
seen, the others can be derived, and let us set for
$\mu_1, \mu_2, \mu_3, \ldots$ the values, arranged in order of magnitude, of the function
\weq{1}{(x-\alpha y)^2 + \beta^2y^2 = n.}

We take $xy$ as rectangular coordinates in the plane
and overlay the plane with two \gesperrt{lattices}, setting in the
one lattice the even integers for the $x$, in the other the odd
ones.

The number of lattice points that lie in the interior of the ellipse (1)
or on its periphery, and for which therefore $\mu_n \leqq n$, we denote for the two kinds by $Z^0(n), Z^1(n)$.

Let us set
\[
x = \xi\sqrt{n}, \qquad y = \eta\sqrt{n},
\]
then (1) goes over into
\weq{2}{(\xi-\alpha\eta)^2 + \beta^2\eta^2 = 1,}
and the area of this ellipse is $\pi/\beta$, and by Bd.\,II,
§ 194, 1.:
%CONT>>
% [sic] The third denominator in(7) is printed2^s, not3^s; retained.
% [sic] The first parenthesized denominator in(8) is printed2, not2^s; retained.
% [sic] In(8) the first epsilon term lacks the factor(1-2^{-s}); printed formula retained.
